\scanpage{27}
% archive sheet number: 26 (inverted)
\mspage{-I- 2 -}
a) \uline{on a} $\underline H'\subset\underline H$ ;

b) \uline{si} $H_{\unsure{o|0}}\in\underline H$ \uline{est tel que} $C'=\displaystyle\bigcap_{H\in\underline H,H\ne H_o}D_H(C)$ \uline{soit distincte de} $C$, \uline{alors} $H_o$ \uline{est un mur de} $C$ ;

c) \uline{on a} $C=\displaystyle\bigcap_{H\in\underline H'}D_H(C)$.

Soient $L\in\underline H'$ un mur de $H$, $x$ un point de $L$ et $U$ un voisinage de $x$ tels que $U\cap C=U\cap L^+$. Si $x$ n'appartenait \unsure{à|a} aucun $H\in\underline H$, il existerait une boule ouverte $B$ de centre $x$ telle que $B\cap H=\varnothing$ pour tout $H\in\underline H$ (car $\underline H$ est localement finie) ; on pourrait supposer que $B\subset U$, de sorte que $B\cap C=B\cap L^+$ ; comme $B\cap C$ est non-vide et que $B$ est connexe, on aurait $B\subset D_H(C)$, de sorte que $B\subset C$, d'où $x\in C$, en contradiction avec $x\in L^+$. Il existe donc $H\in\underline H$ tel que $x\in\underline H$ ; on a $U\cap L^+\subset C\cap D_H(C)$, d'où $H=L$ d'après le lemme 1. On a donc $\underline H'\subset\underline H$, ce qui démontre a).

Démontrons b). On a $C'\not\subset D_{H_o}(C)$. Comme $C'$ est convexe, on a donc $C'\cap H_o\ne\varnothing$ ; notons que $C'\cap H_o$ est relativement ouvert dans $H_o$. Les $H\cap H_o$, pour $H\in\underline H$ et $H\ne H_o$, forment une famille localement finie de sous-espaces affines de $H_o$ ; leur réunion est donc fermée et sans point intérieur, et ne peut par conséquent contenir $C'\cap H_o$. Il existe donc un point $y$ de $C'\cap H_o$ qui n'appartient à aucun $H\in\underline H$ distinct de $H_o$. Comme $\underline H$ est localement finie, il existe un voisinage ouvert convexe $U$ de $y$ tel que $U\cap H=\varnothing$ pour tout $H\in\underline H$, $H\ne H_o$. On a $U\cap C'\ne\varnothing$, d'où $U\cap D_H(C)\ne\varnothing$ pour tout $H\in\underline H$, $H\ne H_o$ ; comme $U$ est un ouvert convexe qui ne rencontre pas $H$, on en déduit $U\subset D_H(C)$. On a donc $U\cap C=U\cap(\displaystyle\bigcap_{H\ne H_o}D_H(C))\cap D_{H_o}(C)=U\cap D_{H_o}(C)$, de sorte que $H_o$ est un mur de $C$. Ceci démontre b).

Passons à c). Soit $z$ un point quelconque de $V$. Joignons-le à un point $y$ de $C$. Comme $\underline H$ est localement finie, le segment $[z,y]$ ne peut rencontrer qu'un nombre fini \margin{\handwritten{$\vert$}}d'hyperplans $H\in\underline H$. Donc les hyperplans $H\in\underline H$ tels que $z\notin D_H(C)$ sont en nombre fini, soient $H_1,\ldots,H_q$. Prenons alors pour $z$ un point de $\displaystyle\bigcap_{H\in\underline H'}D_H(C)$, de sorte que
